El volcà Cumbre Vieja va entrar en erupció el 19 de setembre de 2021 a l'illa de La Palma. El so de l'erupció volcànica va ser mesurat per un equip de vulcanòlegs, que van detectar que estava format per diferents freqüències i que tenia un nivell d'intensitat sonora variant. Pel que fa al so que produeix en alguns moments de l'erupció, un volcà es pot modelitzar com si fos un instrument musical de vent gegant. D'una manera simplificada, podem considerar un volcà com un tub d'aire amb un extrem tancat i l'altre obert, on es produeixen ones estacionàries d'una manera semblant a una trompeta.

\begin{apartat}{1,25}
En un moment determinat de l'erupció del Cumbre Vieja, es va detectar que la freqüència fonamental del so que emetia era de 5,40~Hz. Calculeu la longitud d'ona d'aquest so. Dibuixeu el perfil de l'ona estacionària corresponent i calculeu la longitud del tub, que correspon a la longitud del tram de la xemeneia ple d'aire.
\end{apartat}

\begin{apartat}{1,25}
En un moment de l'erupció, es va detectar un nivell d'intensitat sonora de 72,0~dB a la localitat de Todoque, que és a 5,22~km del con del volcà. Calculeu el nivell d'intensitat sonora que se sentirà des de la costa de l'illa d'El Hierro, que és a 92,3~km del volcà Cumbre Vieja. (Suposeu que no hi ha obstacles en la propagació del so entre el volcà Cumbre Vieja i l'illa d'El Hierro.)
\end{apartat}

\dades{%
  $I_0 = 10^{-12}\ \mathrm{W\,m^{-2}}$.\\
  La velocitat del so en l'aire és de 340~m/s.}
